% 第 12 章（样章）。数字：checks/check_ch12.py；史料引文：data/snapshots/paper_money_facts.csv（逐字核对见
% data/fetch/verify_quotes.py）；元代印钞表：data/snapshots/yuanshi_suyinchao.csv；图：figsrc/fig_ch12.py。
\chapter{宋元明：世界最早的纸币及其崩溃}\label{ch:songming}

\begin{chapterguide}
\guidequestion{纸币为什么最早在中国出现，又为什么一再通胀崩溃？}
\begin{guidegoals}
\item 追踪交子从民间票据到官方纸币（1023 年设益州交子务）的演变，以及它的准备金与分界发行制度；
\item 用会子、元钞与明宝钞的发行和贬值，说明财政压力怎样转化为通胀；
\item 解释明代为什么最终转向白银；
\item 用这一段经验检验“纸币必然失败”与“纸币可以稳定”两种说法。
\end{guidegoals}
\guideprereq{\chref{ch:qinhan}\enspace\nameref{ch:qinhan}；\chref{ch:money}\enspace\nameref{ch:money}}
\end{chapterguide}

今天我们口袋里的钞票，本身几乎一文不值，却能买到东西。\chref{ch:money}说过，这是因为它是中央银行的负债，背后有国家的税收与法律。这种安排第一次大规模出现，不在伦敦或阿姆斯特丹，而在 11 世纪的四川。此后的五百多年里，宋、元、明三个王朝都发行过全国或地区性的纸币，也都经历过它的贬值与崩溃。这一章沿着这条线走一遍：纸币怎样诞生，靠什么维持价值，又是怎样一步步失去价值的。

这段历史的价值在于，它几乎把纸币制度的所有问题都预演了一遍：准备金该留多少，纸币能不能用来纳税，到期不兑换会怎样，战争开支怎样变成印钞，官方定价与市场价格怎样分离。读完这一章，你会发现第~\ref{ch:bullion}、\ref{ch:budget}、\ref{ch:hyperinflation}、\ref{ch:fiscaldominance}~章讨论的很多问题，宋代的官员已经在奏议里争论过了。

\section{当时的世界}\label{sec:songming-world}

\subsubsection{四川的铁钱}
北宋的四川（“川界”）与其他地区不同，通行的是铁钱，而不是铜钱。铁的价值远低于铜，同样的购买力要用重得多的金属。北宋人李攸在《宋朝事实》里记下了这个麻烦：“川界用铁钱，小钱每十贯，重六十五斤，折大钱一贯，重十二斤。”\cite{songchaoshishi} 也就是说，一贯（名义上一千文）小铁钱重 6.5 斤；集市上一笔三五贯的交易，要搬动 19.5—32.5 斤铁。书中接着说：“街市买卖至三五贯文，即难以携持。”

这个困难给了一种替代品机会。成都的富商开始收进铁钱、开出可以随时兑现的纸质凭证，称为\textbf{交子}。持有人兑现时，“每贯割落三十文为利”，即收取 3\% 的手续费。这时的交子还只是一张存款收据：价值完全来自“能换回铁钱”的承诺。承诺一旦动摇，收据就打折扣。私交子后来出现了兑付困难，史书说挤兑时“每一贯多只得七八百”\cite{songchaoshishi}——一千文的收据只能换回七八百文。

\subsubsection{从私交子到官交子}
天圣元年（1023），北宋朝廷接管了交子的发行，在成都设益州交子务；朝廷的命令“天圣元年十一月二十八日到本府”\cite{songchaoshishi}。官交子沿用了私交子的两条规则，又加上了第三条：
\begin{itemize}
\item \textbf{兑现收费}：“每小铁钱一贯文，依例克下三十文入官”，仍是 3\%；
\item \textbf{准备金}（宋人叫\textbf{本钱}）：“大凡旧岁造一界，备本钱三十六万缗”\cite{songshi}；
\item \textbf{分界发行}：“以三年为一界而换之”，“界以百二十五万六千三百四十缗为额”\cite{songshi}。每一界有固定的发行上限，期满用新交子收回旧交子。
\end{itemize}
一缗就是一贯。把两个数字放在一起，每一界 $1\,256\,340$ 缗交子背后有 $360\,000$ 缗现钱，准备率约为 28.7\%（例~\ref{ex:songming-reserve}）。这已经是一张有准备金、有发行上限、有期限的纸币。图~\ref{fig:songming_timeline} 画出了此后五百多年里四种纸币的时间线，以及本章要讨论的事件。

\fig{songming_timeline}{宋元明四种纸币的时间线。看深色条与浅色条的分界：北宋交子在南宋只在四川流通，明宝钞在 15 世纪中叶以后名存实亡。圆点是正文讨论的事件。年代据《宋史》《元史》《明史》的食货志，见本章各节。}{fig:songming_timeline}

\section{一张纸为什么值钱}\label{sec:songming-value}

在讲后来的崩溃之前，先回答一个更基本的问题：一张本身不值钱的纸，凭什么能换东西？从宋元明的经验看，支撑纸币价值的东西有三样。

第一是\textbf{兑现}。如果官府随时按面值把纸币换成铜钱、铁钱或白银，纸币就不会比这些金属便宜太多：一旦打折，人们就会拿去兑现。兑现靠准备金，但准备金不必是百分之百，因为不是所有人同时来兑现。

第二是\textbf{纳税}。如果官府规定某些税可以用纸币缴纳，那么每个纳税人都需要纸币，纸币就有了与兑现无关的需求。南宋的会子一开始就规定，水路不通之处的上供等钱“许尽输会子”\cite{songshi}。

第三是\textbf{数量与期限}。分界发行意味着每一界有上限、到期收回。持有人知道，三年后这张纸会被按面值换成新纸或现钱；官府也无法无限地印下去。用今天的话说，每一界交子都像一笔三年期、不付利息的政府债务。

\begin{intuition}[纸币是不付利息的政府债务]
把交子想成官府向持有人借的钱：持有人交出铁钱，拿到一张欠条；欠条可以转手，到期由官府兑现。只要人们相信官府会按约定兑现、会接受它纳税、不会无限地多写欠条，欠条就能按面值流通。下面各节的每一次崩溃，都是这三条中的一条或几条被打破了。
\end{intuition}

衡量第一条的指标是准备率：用准备金除以发行额。
\begin{keyeq}
\rho=\frac{R}{N},\label{eq:songming-reserve}
\end{keyeq}
其中 $R$ 是本钱（准备金），$N$ 是流通中的纸币面额。$\rho$ 越高，能应付的同时兑现越多；但 $\rho$ 等于 1 时，纸币只是金属的替身，发行者得不到任何好处。宋代的选择是一个中间值。

衡量后两条是否被守住，最直接的办法是看纸币的价值随时间怎样变化。史料往往只给出相隔几十年的两个数字，比较它们时，我们用年均复合变化率：若某个量在 $T$ 年里从 $X_0$ 变为 $X_T$，则
\begin{keyeq}
\bar g=\left(\frac{X_T}{X_0}\right)^{1/T}-1 .\label{eq:songming-avg}
\end{keyeq}
它回答的是：如果每年按同样的比例变化，要多大的比例才能在 $T$ 年里累积出这样的倍数。它掩盖了中间的起伏，所以只用来比较不同时期的平均速度。本章所有“年均”数字都用式~\eqref{eq:songming-avg} 计算。

\begin{example}{交子的准备率与兑现费}\label{ex:songming-reserve}
天圣年间每一界交子以 $1\,256\,340$ 缗为额，备本钱 $360\,000$ 缗；兑现时每贯收 30 文。求准备率，以及兑现费占面值的比例。若某一年持有人要求兑现的交子超过发行额的多少，本钱就会用完？
\solution 由式~\eqref{eq:songming-reserve}，$\rho=360\,000/1\,256\,340\approx 28.7\%$。一贯为 1000 文，兑现费为 $30/1000=3\%$。只要同时来兑现的不超过发行额的 28.7\%，本钱就够用；超过这个比例，官府就必须另找现钱，否则只能拒兑。3\% 的兑现费是交子务的收入，也让人们不愿意为小事兑现，从而减少了对本钱的消耗。
\end{example}

\begin{compare}{可兑换纸币与不可兑换纸币}{可兑换纸币}{不可兑换纸币}
\cmprow{价值的锚}{承诺按面值换成金属；价值不会远离兑换比价}{只有纳税需求与对发行数量的预期}
\cmprow{约束发行的}{准备金：多发就要多备本钱，否则承受挤兑}{只有发行者的自我约束}
\cmprow{出问题时}{先表现为挤兑与停兑}{先表现为物价缓慢上涨，往往持续多年}
\cmprow{本章的例子}{北宋交子、元初的中统钞（1260—1275）}{明宝钞；元 1310 年以后的纸币}
\cmprow{今天的对应}{金本位、货币局（\chref{ch:emcrises}）}{各国现行的法定货币（\chref{ch:money}）}
\cmpsum{可兑换纸币的价值由准备金守住，不可兑换纸币的价值只能由发行者的财政纪律守住。}
\end{compare}

\begin{compare}{分界发行与无期限发行}{分界发行}{无期限发行}
\cmprow{规则}{每一界有定额，三年期满以新换旧}{一经发行永久流通，旧钞只换新纸不收回}
\cmprow{像什么}{三年期、不付利息的政府债务}{永续、不付利息的政府债务}
\cmprow{违约的样子}{“展界”：到期不兑，延期再用}{无所谓到期，只能看物价}
\cmprow{本章的例子}{交子、早期会子}{1247 年以后的会子、元钞、明宝钞}
\cmpsum{分界让每一次超发都暴露在到期日上；取消分界，就取消了这个到期日的约束。}
\end{compare}

\section{北宋：从交子到钱引}\label{sec:songming-north}

官交子的第一个破绽来自边防开支。西北用兵时，朝廷曾从四川调用交子。皇祐三年（1051），田况追述此事，说朝廷“兼自秦州两次借却交子六十万贯，并无见钱桩管”\cite{songchaoshishi}：六十万贯交子发出去了，却没有相应的现钱作本钱。这是用纸币直接为财政融资的较早记录。与例~\ref{ex:songming-reserve} 对照：一界的本钱只有三十六万缗，这一笔借用就超过了它。

到徽宗大观元年（1107），问题已经无法掩盖。《宋史·食货志》的概括是“大观中，不蓄本钱而增造无艺”：当年发行“较天圣一界逾二十倍，而价愈损”\cite{songshi}。朝廷的对策是换发新交子，规定“新交子一当旧者四”：一张旧交子只能换四分之一张新交子。对持有人来说，这等于官府赖掉了 75\% 的债。次年张持报告，四川“本路引一千者今仅直十之一”，面值一千的钱引只值一百，损失 90\%；史书还说，当时“至引一缗当钱十数”\cite{songshi}。

但北宋的故事里也有一个反例。张商英执政后，“奉诏复循旧法”，恢复旧日的定额与分界；宣和年间他奏报说，“自旧法之用，至今引价复平”\cite{songshi}。这是当事人自己的说法，未必没有夸大，但它说明了宋人的理解：同一种纸币，限量、分界的规则一旦恢复，价值就能回到面值附近。纸币的贬值并不是无可挽回的物理过程，而是取决于规则是否被守住。

北宋亡后，四川的纸币（此时叫钱引）在南宋继续流通，规则却越来越松。到绍兴末年（1162 年前后），四川钱引“积至四千一百四十七万余贯；所贮铁钱，仅及七十万贯，以盐酒等阴为称提”\cite{songshi}。流通额是北宋一界定额的约 33 倍，铁钱准备只有流通额的 1.7\%。史文用了一个关键的词：准备不足的部分，靠盐、酒专卖的收入“阴为称提”。现钱准备撤退之后，接替它的是财政收入。

\section{南宋：会子与称提}\label{sec:songming-huizi}

绍兴三十年（1160），户部侍郎钱端礼“被旨造会子，储见钱，于城内外流转”\cite{songshi}：最初的\textbf{会子}同样有现钱储备。南宋的纸币从此走出四川，先行于两浙，后来通行于淮、浙、湖北、京西。会子的设计吸取了交子的经验：乾道四年（1168）定为“三年立为一界，界以一千万贯为额”，每道收“靡费钱二十足”作工本\cite{songshi}。前面说过，会子从一开始就能用来缴纳上供等钱。

会子在财政里的分量很快变得很大。淳熙三年（1176），“户部岁入一千二百万，其半为会子”\cite{songshi}：每年 1200 万贯的收入中，600 万贯是会子。官府收进会子，再把它花出去，会子就在财政收支中循环。

\subsubsection{称提：南宋的公开市场操作}
宋人把维持纸币价值的种种办法称为\textbf{称提}。《宋史》追述高宗论四川交子时，“最善沈该称提之说”，即“官中常有钱百万缗，如交子价减”，就由官府出钱收买\cite{songshi}。乾道二年（1166），“以会子之弊”，朝廷“出内库及南库银一百万收之”\cite{songshi}。这正是今天中央银行公开市场操作的逻辑：用资产（银、钱）买回负债（纸币），减少流通量，托住价格。

\begin{compare}{称提与公开市场操作}{南宋的称提}{今天的公开市场操作}
\cmprow{做法}{出金银、现钱、度牒收买会子；以旧换新时打折收回}{央行买卖国债，改变银行体系的准备金}
\cmprow{目标}{会子的价格（对铜钱的比价）}{短期利率或物价（第~\ref{ch:money}、\ref{ch:consensus}~章）}
\cmprow{资产从哪来}{内库与封桩库的金银、专卖收入}{央行资产负债表，最终是财政}
\cmprow{失灵的条件}{收回赶不上新印；用强制摊派代替买回}{财政主导时，央行不得不为赤字融资（\chref{ch:fiscaldominance}）}
\cmpsum{两者都是用资产买回货币以维持其价值；能做多少，最终取决于背后的财政。}
\end{compare}

\subsubsection{展界：到期不兑}
分界的约束在南宋逐渐松动。绍熙元年（1190），第七、第八界会子到期后又各展三年，有臣僚抗议：“会子界以三年为限，今展至再，则为九年，何以示信？”\cite{songshi} 到期不兑而延期，叫\textbf{展界}。用债务的语言说，展界就是政府无力偿还到期债务时的强制展期。在四川，“自开禧军兴以后，用度不给，展年收兑，遂至两界、三界通使”；在两淮，交子“其数日增，价亦日损，称提无术，但屡与展界而已”\cite{songshi}。战争开支、展界、几界同时流通、价格下跌，这几件事总是一起出现。

表~\ref{tab:songming-huizi} 把《宋史·食货志》中会子的几个数量排在一起。定额从 1168 年的一千万贯提高到 1195 年的三千万贯；到嘉定二年（1209），仅第十二、十三两界尚在流通的就有一亿零二百余万贯，是 1168 年一界定额的 10.2 倍。这一年朝廷拨出黄金一百五万两、每两作价四十贯，即 4200 万贯，加上度牒、官告等收兑旧会；黄金只抵得上两界流通额的 41\%，于是规定“以旧会之二，易新会之一”\cite{songshi}——旧会子持有人一次损失一半。史书对这一年的概括是：“以三界会子数多，称提无策。”

\begin{table}[htbp]\centering\small
\caption{南宋会子的定额与流通额（《宋史·食货志下三》）}\label{tab:songming-huizi}
\begin{tabularx}{\linewidth}{@{}l Y r r@{}}
\toprule
年份 & 史文 & 数量（贯） & 相当于 1168 年一界定额 \\\midrule
1168 & 三年立为一界，界以一千万贯为额 & 1000 万 & 1 \\
1176 & 户部岁入一千二百万，其半为会子 & 600 万（岁入中的会子） & — \\
1195 & 诏会子界以三千万为额 & 3000 万 & 3 \\
1209 & 十二界、十三界除烧毁尚有一万二百余万贯 & 1.02 亿 & 10.2 \\
1232 & 两界会子已及三亿二千九百余万 & 3.29 亿 & 32.9 \\
     & 异文：二亿二千九百余万 & 2.29 亿 & 22.9 \\
\bottomrule
\end{tabularx}
\datasource{《宋史》卷一百八十一，殆知阁数字文本；1232 年异文据维基文库本。数字整理与比值见 \texttt{data/snapshots/paper\_money\_facts.csv} 与 \texttt{checks/check\_ch12.py}。}
\end{table}

绍定五年（1232）的数字有异文：我们使用的数字文本与《钦定续文献通考》作“三亿二千九百余万”，维基文库本《宋史》作“二亿二千九百余万”。两种读法都保留，例~\ref{ex:songming-huizi} 会说明这种不确定性对结论的影响其实不大。

\begin{example}{会子的增长与异文}\label{ex:songming-huizi}
1168 年会子一界定额为 1000 万贯；1232 年两界流通额为 3.29 亿贯（异文 2.29 亿贯）。分别求两种读法下流通额相对 1168 年定额的倍数，以及 64 年里的年均增长率。异文对结论有多大影响？
\solution 倍数为 $3.29\text{ 亿}/1000\text{ 万}=32.9$ 与 $22.9$。由式~\eqref{eq:songming-avg}，年均增长率为 $32.9^{1/64}-1\approx 5.6\%$ 与 $22.9^{1/64}-1\approx 5.0\%$。两种读法的水平相差约 44\%，年均增长率却只差 0.6 个百分点：在几十年的时间上开方，水平上的不确定性被大大压缩了。无论取哪一种，结论都是会子以每年 5\%—6\% 的速度持续增加，远远超过了分界制度设想的“每界定额、到期收回”。
\end{example}

结局是可以预料的。淳祐七年（1247），朝廷宣布“以十八界与十七界会子更不立限，永远行使”\cite{songshi}：分界制度正式取消，会子成了无期限的纸币。咸淳四年（1268），朝廷又发行“见钱关子”，规定关子一贯作 770 文足，第十八界会子“每道作二百五十七文足”\cite{songshi}。一道面值一贯的会子，官定只值关子的 33.4\%，铜钱面值的 25.7\%。八年后临安降元。

\section{元：以银为本的全国纸币}\label{sec:songming-yuan}

元朝第一次把纸币作为全国通行的主要货币\cite{guan2024}。中统元年（1260）发行\textbf{中统钞}，规定“每一贯同交钞一两，两贯同白银一两”\cite{yuanshi}：两贯中统钞兑一两白银。《元史·食货志》用一句话概括它的原理：“其法以物为母，钞为子，子母相权而行。”\cite{yuanshi} 中统五年（1264）设各路平准库以平物价，“仍给钞一万二千锭，以为钞本”\cite{yuanshi}。锭是元代纸币的大单位，一锭为五十贯（本章的“深入一步”）。

Guan、Palma 与 Wu（2024）\cite{guan2024}把元代纸币分为三个阶段：1260—1275 年白银可以兑换，银本位真正约束着发行；1276—1309 年名义上可兑换，实际上兑换处有限、要打折扣；1310 年起成为不兑现的纸币。下面的数字与这个分期是一致的。

\subsubsection{发行：一张表里的七十年}
《元史》卷九十三保存了一张罕见的表：从中统元年到天历二年（1260—1329），每一年印了多少锭中统钞、至元钞与至大银钞\cite{yuanshi}。七十年里每年都有记录。把三种钞按官定比价折成中统钞（至元钞一贯当中统钞五贯，至大银钞一两当至元钞五贯），就得到图~\ref{fig:yuan_issuance}。

\fig{yuan_issuance}{元代每年印钞数（对数纵轴，折合中统钞）。看 1276 年的跳升与 1310 年的尖峰：前者是元军攻下临安的那一年，后者是发行至大银钞的那一年。虚线是 1260—1275 年与 1276—1286 年两个时期的年均印钞数；1287 年以后的至元钞与 1310 年的至大银钞已按官定比价折算。数据：《元史》卷九十三“岁印钞数”（\texttt{data/snapshots/yuanshi\_suyinchao.csv}），与 Guan 等（2024）附表 A1 中 1260、1264、1265、1276、1281、1287 年的数字一致。}{fig:yuan_issuance}

图中最醒目的是 1276 年：这一年元军攻下临安，南宋事实上灭亡，印钞从前一年的 39.8 万锭跳到 142 万锭。1260—1275 年的 16 年里，每年平均印 10.6 万锭；1276—1286 年的 11 年里，每年平均印 117.4 万锭，是前一时期的 11.1 倍。准备金也远远跟不上：1264 年拨给平准库的一万二千锭钞本，只相当于 1260—1264 年累计印钞（35.6 万锭）的 3.4\%。

\subsubsection{贬值：从官定比价里读出来}
至元二十四年（1287），元朝改发\textbf{至元钞}，“每一贯文当中统钞五贯文”；同时规定花银每两“入库其价至元钞二贯”，赤金每两“入库二十贯”\cite{yuanshi}。官府自己定的这些比价，透露了中统钞贬值了多少（例~\ref{ex:songming-yuan}）：1260 年一两白银值中统钞二贯，1287 年值至元钞二贯，也就是中统钞十贯。按官定银价，中统钞只剩发行时价值的 20\%。《元史》说这一套办法“其法为最善”\cite{yuanshi}，但这一年按中统钞计的印钞量达到 508.8 万锭，1287—1290 年平均每年 527.5 万锭。

二十二年后，同样的事情又发生了一次。至大二年（1309），武宗发行至大银钞，诏书承认至元钞行用“逮今又复二十三年，物重钞轻，不能无弊”，规定至大银钞“每一两准至元钞五贯，白银一两，赤金一钱”\cite{yuanshi}。一两白银从至元钞二贯涨到五贯，至元钞按官定银价只剩 40\%。《元史》的总结是：“大抵至元钞五倍于中统，至大钞又五倍于至元。”\cite{yuanshi} 1310 年按中统钞计的印钞量达到 3630 万锭，是 1260 年的 494 倍。这个尖峰没有持续：至大四年（1311），仁宗即位后下诏，说至大银钞“钞以倍数太多，轻重失宜”，予以废除，恢复中统、至元钞\cite{yuanshi}。

\begin{example}{元钞的官定贬值}\label{ex:songming-yuan}
1260 年：二贯中统钞兑白银一两。1287 年：白银一两入库价为至元钞二贯，至元钞一贯当中统钞五贯。1309 年：白银一两准至元钞五贯。求 (1) 1287 年中统钞相对 1260 年的价值与银价的年均涨幅；(2) 1309 年至元钞相对 1287 年的价值与年均涨幅；(3) 1260—1309 年，按中统钞计的银价上涨了多少倍。
\solution (1) 1287 年白银一两值至元钞二贯，即中统钞 $2\times5=10$ 贯；中统钞的银价值为 $2/10=20\%$。由式~\eqref{eq:songming-avg}，银价年均涨幅为 $(10/2)^{1/27}-1\approx6.1\%$。
(2) 白银一两从至元钞二贯涨到五贯，至元钞价值为 $2/5=40\%$，年均涨幅 $(5/2)^{1/22}-1\approx4.3\%$。
(3) 1309 年白银一两值至元钞五贯，即中统钞 25 贯；与 1260 年的二贯相比，上涨 $25/2=12.5$ 倍。

这些都是\emph{官定}比价，不是市场成交价；它们说明的是官府自己承认的贬值幅度。1287 年官定的金银比价为 $20/2=10$。
\end{example}

\begin{pitfall}[官定比价就是市场价格，盐价就是物价]
官定比价往往落后于市场：官府只在无法维持时才承认贬值，市场价格通常更早、更多地下跌。另一方面，单一商品的价格也不等于物价水平。元统二年（1334）御史台的奏报说，大德年间中统钞一贯可买盐四斤八两（4.5 斤），到泰定年间“钞一贯，仅买盐一斤”\cite{yuanshi}，盐价涨了 4.5 倍；但奏报把这次上涨归因于罢除官府卖盐的局所、恢复商人贩卖，即盐的流通制度的变化，而不是纸币。Guan 等人用多种价格指数做的估计表明，元代物价在 1290 年代到 1330 年代大体平稳\cite{guan2024}。判断纸币是否贬值，要看多种商品与金银的价格，并分清货币原因与商品本身的原因。
\end{pitfall}

\subsubsection{至正：最后的崩溃}
元末的崩溃来得很快。至正十一年（1351）改发至正交钞，正值各地起事；《元史》说此后“物价腾踊，价逾十倍”，到后来“京师料钞十锭，易斗粟不可得”\cite{yuanshi}——十锭即五百贯，换不到一斗粮食。Guan 等人的研究也发现，元代的超发主要由战争开支、特别是内战驱动，而且在不兑现的时期，战争引起的超发要大得多\cite{guan2024}。

\begin{case}{伊利汗国 · 1294 年 · 两个月的纸币}
元朝的纸币制度曾被搬到波斯。伊利汗国的乞合都汗（Gaykhatu，1291—1295 年在位）为解决财政困难，仿照元钞发行纸币（波斯语称 čāv，即“钞”）。据史家拉施特（Rashid al-Din, 1247—1318）的记载，纸币于 1294 年 9 月 12 日在大不里士开始流通，同时禁止用金属货币交易。结果物价上涨了十倍以上，许多居民在几天之内逃离城市，市场停摆。混乱持续了约两个月，官府先是允许用铸币买粮，最后撤销了纸币令\cite{jackson1990}。

同样是纸币，为什么在元朝能流通几十年，在大不里士只撑了两个月？对照\secref{sec:songming-value}的三个条件：伊利汗国的纸币没有准备金，不能纳税换来信任，也没有任何发行纪律的记录；它唯一的支撑是禁止使用金银的命令。仅靠强制，纸币一天也撑不住。如果十倍的上涨发生在两个月之内，平均每月的涨幅就超过 216\%，远远高于本章最后一节要用到的恶性通胀门槛。
\end{case}

\section{明：没有本钱、没有分界的宝钞}\label{sec:songming-ming}

明朝的纸币从一开始就缺少前两个朝代曾有过的约束。洪武八年（1375）发行\textbf{大明宝钞}，“每钞一贯，准钱千文，银一两；四贯准黄金一两”，同时“禁民间不得以金银物货交易，违者罪之；以金银易钞者听”\cite{mingshi}：可以用金银换钞，不能用钞换金银。宝钞既不兑现，也不分界，没有像宋代本钱那样的准备金\cite{vonglahn1996,pengxinwei}。支撑它的只有两样东西：禁止金银交易的法令，以及部分税收收钞——洪武十年（1377）规定商税“钱三钞七”，七成收钞\cite{mingshi}。

后果来得很快。洪武十八年（1385），官员俸禄折钞，“二贯五百文准米一石”；到洪武二十六年（1393）前后，两浙、江西、福建、广东已经“有以钱百六十文折钞一贯者”\cite{mingshi}，即按铜钱计，宝钞只剩面值的 16\%。洪武三十年（1397）朝廷再次“申交易用金银之禁”，永乐初年又加重刑罚，仍然无济于事。

明代的官员对原因看得很清楚。永乐二年（1404），都御史陈瑛说：“比岁钞法不通，皆缘朝廷出钞太多，收敛无法，以致物重钞轻。”他建议按人口纳钞领盐，“诚令计口纳钞食盐，可收五千馀万锭”，用税收把钞收回来\cite{mingshi}。仁宗即位（1424）时问户部尚书夏原吉钞法为何不行，夏原吉答：“钞多则轻，少则重。民间钞不行，缘散多敛少，宜为法敛之。”他的对策有两条：加征可以用钞缴纳的门摊等商税，把钞收回官府；以后“自今官钞宜少出，民间得钞难，则自然重矣”\cite{mingshi}。“钞多则轻，少则重”正是货币数量论的核心命题，比博丹（Jean Bodin, 1530—1596；\chref{ch:silver}）与休谟（David Hume, 1711—1776；\chref{ch:bullion}）的论述都早。南宋的黄畴若在嘉定初年也说过：“物少则贵，多则贱，理之常也。”\cite{songshi}

但知道原因不等于能够收手，收手了也不一定见效。史书说夏原吉的办法施行之后，“然是时，民卒轻钞”\cite{mingshi}。宣德初年（1426），“米一石用钞五十贯”，是 1385 年俸禄折价的 20 倍；成化年间（1465—1487），“钞一贯不能直钱一文”；嘉靖四年（1525），“钞一贯折银三厘”，即面值一两白银的宝钞只值 0.003 两\cite{mingshi}。图~\ref{fig:note_value} 把元、明两代纸币的价值画在同一种刻度上。

\fig{note_value}{纸币价值相对发行时的变化（对数纵轴，发行时 = 100，纵轴每一格是十倍）。看两条线的斜率：斜率就是年均贬值速度。(a) 元：按官定银价，中统钞到 1287 年剩 20，至元钞到 1309 年剩 40。(b) 明：按铜钱计，宝钞到 1393 年前后剩 16，成化年间（图中画在 1475 年）不足 0.1；按白银计，1525 年剩 0.3。数据：《元史》卷九十三、《明史》卷八十一（\texttt{data/snapshots/paper\_money\_facts.csv}）。}{fig:note_value}

\begin{example}{明宝钞的贬值速度}\label{ex:songming-ming}
(1) 1385 年俸禄折价为米一石二贯五百文，1426 年米一石用钞五十贯。求米价（以钞计）上涨的倍数与年均涨幅。(2) 1375 年宝钞一贯准银一两，1525 年一贯折银三厘（0.003 两）。求以白银计的宝钞价值与银价（以钞计）的年均涨幅。(3) 与 1393 年前后的铜钱折价比较，哪一段贬值最快？
\solution (1) $50/2.5=20$ 倍；年均涨幅 $20^{1/41}-1\approx7.6\%$。前者是官定的俸禄折价，后者是市场米价，口径不完全相同，结果只表示数量级。
(2) 宝钞价值为 $0.003/1=0.3\%$；银价年均涨幅为 $(1/0.003)^{1/150}-1\approx3.9\%$。
(3) 1375—1393 年，按铜钱计的钞价从一千文跌到一百六十文，铜钱的钞价年均上涨 $(1000/160)^{1/18}-1\approx10.7\%$，是三段中最快的。150 年平均每年 3.9\%，看起来不高，累积起来却让宝钞失去了 99.7\% 的价值。
\end{example}

宝钞的失败让明代的经济转向了白银。民间交易越来越多地用银，赋役也逐步折银征收；到隆庆年间，史书已经说“宝钞不用垂百馀年，课程亦鲜有收钞者，惟俸钱独支钞如故”\cite{mingshi}：税收早已不收钞，只有官员的俸禄还用这种不值钱的纸来发。一个没有自己能控制的货币、只能依靠进口白银的帝国，就这样把货币供给交给了海外的银矿与马尼拉的帆船\cite{vonglahn1996}。这是\chref{ch:silver}的故事。

\section{数据检验与反例}\label{sec:songming-test}

表~\ref{tab:songming-systems} 把四种纸币放在一起。比较它们，可以检验本章开头的两种说法。

\begin{table}[htbp]\centering\small
\caption{宋元明四种纸币的制度与结局}\label{tab:songming-systems}
\begin{tabularx}{\linewidth}{@{}l Y Y Y Y@{}}
\toprule
 & 北宋交子 & 南宋会子 & 元钞 & 明宝钞 \\\midrule
首次发行 & 1023 年，益州交子务 & 1160 年，临安 & 1260 年，中统钞 & 1375 年 \\
兑现与准备 & 本钱约为一界定额的 28.7\% & 初“储见钱”；后靠金银收兑（称提），无固定比例 & 1260—1275 年可兑银；平准库 & 不兑现，无本钱 \\
期限 & 三年一界 & 三年一界，屡次展界；1247 年起无期限 & 不分界 & 不分界 \\
税收 & — & 可纳上供等钱；岁入一半为会子（1176） & — & 商税七成收钞（1377） \\
贬值的记录 & 1107 年超发二十倍，新一当旧四 & 1268 年每道官定 257 文 & 官定银价：1287 年剩 20\%，1309 年至元钞剩 40\% & 1525 年一贯折银三厘 \\
\bottomrule
\end{tabularx}
\datasource{本章各节所引《宋朝事实》《宋史》《元史》《明史》；元代可兑换的分期据 Guan 等（2024）。}
\end{table}

\subsubsection{“纸币必然失败”？}
四种纸币最终都失去了价值，看起来支持这种说法。但细看每一次失败，原因都不是“纸”本身，而是某一条约束被打破：北宋是边防借用与大观年间的超发，南宋是战争、展界与取消分界，元是灭宋与内战时期的印钞，明是从一开始就没有准备与期限。反过来，约束被守住的时候，纸币是能稳定的：大观之后张商英恢复旧法，据他奏报“引价复平”；元代物价在 1290 年代到 1330 年代大体平稳\cite{guan2024}。所以更准确的说法是：纸币的价值取决于发行者能否守住约束，而在宋元明的条件下，战争一来，约束几乎总是先被放弃。

\subsubsection{两种解释}
纸币为什么贬值，宋元明的材料支持两种不同侧重的解释，把它们并排放在一起比较。

\begin{compare}{数量论解释与背书解释}{数量论解释}{背书（财政）解释}
\cmprow{核心命题}{“钞多则轻，少则重”：价值取决于流通数量}{价值取决于发行者以后能否用资产或税收收回纸币}
\cmprow{本章的证据}{会子年增 5\%—6\%；元 1276 年以后年均印钞增至 11 倍，官定银价随后上涨}{钱引“以盐酒等阴为称提”；会子可以纳税；伊利汗国无背书即崩溃}
\cmprow{解释不了的}{为什么同样的数量，有人信、有人不信}{背书的大小很难直接观测}
\cmprow{后面的章节}{\chref{ch:bullion}（数量论），\chref{ch:hyperinflation}（卡甘模型）}{\chref{ch:budget}（不愉快的算术），\chref{ch:fiscaldominance}（财政价格理论）}
\cmpsum{数量论说多印会贬值；背书解释说，人们对未来能否收回的预期决定了多印多少才会贬值。}
\end{compare}

两种解释并不矛盾。陈瑛的话里两者都有：“出钞太多”是数量，“收敛无法”是背书。夏原吉的对策也是两条并用：以后少发，同时用商税把钞收回来。可是在宝钞不兑现、官府还要靠它支付俸禄赏赐的情况下，人们并不相信少发会长久，结果“民卒轻钞”。称提的成败、展界的后果，都说明人们关心的不只是今天有多少纸币，还有官府将来能否、愿不愿意把它们收回去。

\subsubsection{慢性通胀，而不是恶性通胀}
卡甘（Phillip Cagan, 1927—2012）在研究 20 世纪的恶性通胀时，把物价每月上涨 50\% 定为恶性通胀的门槛\cite{cagan1956}。每月 50\% 相当于每年上涨 $1.5^{12}\approx129.7$ 倍。用式~\eqref{eq:songming-avg} 算出的本章各段的年均速度，大约在每年 4\%—11\% 之间：会子的流通额每年增加 5\%—6\%，中统钞的银价每年上涨 6.1\%，明宝钞的铜钱折价每年 10.7\%。除了 1294 年的大不里士与元末的短暂时期，宋元明的纸币贬值大多是\textbf{慢性}的：每年看起来不多，持续几十年、上百年，复利累积下来却几乎吃掉全部价值。魏玛（\chref{ch:weimar}）与\chref{ch:hyperinflation}的恶性通胀是另一回事：它们在一两年内完成同样的过程。慢性通胀给了官府一次又一次换发新钞、宣布新比价的时间，也让每一次换发都成了对持有人的一次部分违约。

\section{今天还剩什么}\label{sec:songming-legacy}

宋元明的纸币早已消失，但它们留下的问题全部还在。

\textbf{准备金与兑现。}交子的本钱、元初的平准库，今天对应金本位下的黄金准备与货币局制度（\chref{ch:emcrises}）：用准备金守住纸币的价值，代价是放弃印钞的自由。准备率越低，越依赖人们不同时来兑现。

\textbf{纸币是政府的负债。}分界发行把这一点写在了纸面上：每一界是一笔三年期的债务，展界就是强制展期，“新交子一当旧者四”与“以旧会之二，易新会之一”就是部分违约。今天的法定货币没有到期日，但它仍然是负债，它的价值仍然取决于政府的预算约束（\chref{ch:budget}）与财政是否主导货币（\chref{ch:fiscaldominance}）。

\textbf{私人发行的“交子”。}成都富商的私交子是存款收据，收兑现费，在挤兑时“每一贯多只得七八百”。今天的稳定币也承诺按面值兑回法定货币，也依赖发行者的准备资产；它们会遇到什么问题，\chref{ch:digitalmoney}会讨论。

\textbf{约束要靠制度，不能只靠认识。}从黄畴若到夏原吉，宋明官员对“多则贱”的道理并不陌生。纸币仍然一再崩溃，是因为在战争与财政压力面前，没有一个能够拒绝印钞的制度。现代的中央银行独立性与通胀目标制（第~\ref{ch:credibility}、\ref{ch:consensus}~章），可以看作对这个问题的一种回答；它是否足够，还要看财政主导的压力有多大。

\begin{deeper}[单位、古代数字与异文]
\textbf{单位。}一贯（一缗）名义上为一千文。南宋常用“省陌”：一贯实付七百七十文（“文足”指足数的文）；咸淳四年关子“贯作七百七十文足”即是此意。元代纸币以锭计，《钦定续文献通考》引孙承泽语：“故五十贯为一锭。”\cite{xuwenxiantongkao} 明宝钞一贯名义上准钱千文、银一两。

\textbf{古代数字。}古文数字中间的零往往不写，例如“百二十五万六千三百四十”是 $1\,256\,340$，“一百九万四千八百”是 $1\,094\,800$ 而不是 $1\,904\,800$；首位的“一”也可以省去，如“百二十五万”。通用的中文数字转换工具按现代写法读数，会把这类数字读错。本书用一个小的读数程序（\texttt{data/fetch/classical\_numerals.py}）逐条读出，并把读出的数字按古文写法写回，与原文逐字比对，两者一致才采用。

\textbf{异文。}史书在流传中有异文。本章的引文据殆知阁古代文献的数字文本\cite{songshi}，并与维基文库本《宋史》卷一百八十一逐条比对：除绍定五年（1232）的“三亿”与“二亿”之外均一致，两种读法都保留（表~\ref{tab:songming-huizi}）。元代印钞数与 Guan 等（2024）的整理一致；明代数字只有《明史》一个来源。
\end{deeper}

\begin{takeaways}
\begin{itemize}
\item 纸币最早出现在北宋四川，起因是铁钱太重；1023 年官府接管交子，规定了兑现费（3\%）、本钱（约为一界定额的 28.7\%）与三年一界的定额发行。
\item 支撑纸币价值的有三样：兑现、纳税与发行数量的约束。宋元明的每一次崩溃，都是其中一条或几条被打破，而打破它们的通常是战争开支。
\item 分界发行使纸币像一笔不付利息的短期政府债务；展界是强制展期，“以旧换新”时的折价是部分违约。
\item 宋元明的纸币贬值大多是慢性的：年均几个百分点，持续几十年至上百年，最终几乎丧失全部价值；这与 20 世纪一两年内完成的恶性通胀不同。
\item 宝钞失败后，明代转向白银，把货币供给交给了海外银矿，这是\chref{ch:silver}的起点。
\end{itemize}
\end{takeaways}

\begin{selfcheck}
\begin{enumerate}
\item 大观元年“新交子一当旧者四”，嘉定二年“以旧会之二，易新会之一”。分别说明旧纸币持有人损失了多少，为什么说这是部分违约。
\item 南宋会子能够纳税，为什么仍然持续贬值？用\secref{sec:songming-value}的三个条件说明。
\item 假定 1168 年的会子按面值流通，到 1268 年第十八界会子每道官定为 257 文足。求这一百年里以铜钱计的会子价值的年均下降速度对应的铜钱价格年均涨幅，并评论“年均很小”与“百年后所剩无几”的关系。
\item 元 1287 年官定金银比价为 10，明 1375 年“四贯准黄金一两”而一贯准银一两。两者相差很大，可以从中推出明初金银的市场比价吗？为什么？
\item 为什么伊利汗国的纸币只维持了两个月，而元钞流通了几十年？
\item 夏原吉说“钞多则轻，少则重”，主张“为法敛之”。如果官府真的减少了宝钞流通量，宝钞的价值一定会回升吗？什么条件下不会？
\end{enumerate}
\end{selfcheck}

\begin{answers}
\item 旧交子一张只值新交子四分之一，损失 75\%；旧会子两道换新会子一道，损失 50\%。官府发行的纸币是它的负债，以低于面值的比率换回，就是只偿还了一部分债务。
\item 纳税只创造了有限的需求：每年的税收只能吸收一部分会子。当发行（年均增长 5\%—6\%）远远超过税收能够收回的数量，又不断展界、最后取消分界时，人们预期会子会越来越多，纳税需求就不足以维持面值。
\item 由式~\eqref{eq:songming-avg}，铜钱价格（以会子计）的年均涨幅为 $(1000/257)^{1/100}-1\approx1.4\%$。每年只有 1.4\%，一百年后会子却只剩四分之一的价值：慢性通胀的危害在于复利与时间。（1168 年会子是否足值，史料并未给出，这里只是假定。）
\item 不能。两者都是官府规定的比价，不是市场成交价，也不一定是同一时期的比较。它们只说明官府在各自的时间给金和银定了怎样的价格。至于中国白银相对黄金较贵、因而吸引白银流入的问题，要用市场价格来讨论（\chref{ch:silver}）。
\item 元钞初期有白银兑换与平准库，可以纳税，发行受银本位约束；伊利汗国的纸币没有准备、没有纳税的用途，也没有发行纪律可言，唯一的支撑是禁止使用金银的命令，一旦市场停摆，命令也只能撤销。
\item 不一定。如果人们相信官府以后还会大量印钞（例如战争仍在继续、俸禄赏赐仍要用钞支付），今天的回收只会被看作暂时的，价值未必回升。夏原吉的办法施行之后“民卒轻钞”，两年后米价仍是 1385 年的 20 倍，就是这种情形。只有当减少发行与收回被看作长期的规则，即有可信的背书时，数量的减少才会转化为价值的回升。第~\ref{ch:credibility}、\ref{ch:hyperinflation}~章会回到这个问题。
\end{answers}

下一章从明代转向白银说起：大量美洲白银怎样经马尼拉流入中国，又怎样在欧洲引起了一场“价格革命”。
